% Source: sources/returns/S6_04_MAPS_MANIFOLDS/sections/manifolds.tex:615–682
\paragraph{Lemma D.1.10 (Proved: the cubic detects the complete fundamental-class invariant).}
Let $A=\mathbb Z^r\oplus\mathbb Z/d$ with $\gcd(d,6)=1$.
The integral homology group of $K(A,2)$ in degree six is
\[
 H_6(K(A,2);\mathbb Z)\simeq
 \mathbb Z^{\binom{r+2}{3}}\oplus
 (\mathbb Z/d)^{\binom{r+1}{2}+r+1}.
\]
Its free coordinates are detected by integral degree-two cubics;
its torsion coordinates are detected by the finite-coefficient
monomials $u x_i x_j$ ($i\leq j$), $u^2x_i$ and $u^3$, where
$u$ evaluates to one on the cyclic summand of $A$. Consequently
D.1.9 and the integral cubic determine the complete class
$\mu(\widetilde X_p)$ in Zhubr's theorem, not merely a quotient of it.
\begin{proof}
We give the low-degree calculation to exclude invisible torsion.
Put $R=\mathbb Z[1/6]$. In the path-loop Serre spectral sequence of
$K(\mathbb Z,3)$ the fibre is $K(\mathbb Z,2)=\mathbb CP^\infty$.
If $t$ has degree two, its transgression is the degree-three
generator $z$. Multiplicativity gives
$d_3(t^2)=2tz$ and $d_3(t^3)=3t^2z$.
Since two and three are units and the total path space is
contractible, induction in total degree forces the base cohomology
through degree seven to be the exterior algebra on $z$:
there is no surviving fibre class that could kill an additional
base class in degrees four through seven. Also $z^2=0$ because
$z$ is odd and two is a unit. This proves exactly the low-degree
base calculation needed below.

The coefficient extension
$0\to\mathbb Z\xrightarrow{d}\mathbb Z\to\mathbb Z/d\to0$
gives the fibration
\[
 K(\mathbb Z,2)\longrightarrow K(\mathbb Z/d,2)
 \longrightarrow K(\mathbb Z,3)
\]
with transgression $d_3(t)=dz$. Its cohomology through degree seven
is therefore computed by $R[t]\otimes\Lambda(z)$ with this
 differential. The groups in degrees three, five and seven are
$R/(d)$, $R/(2d)$ and $R/(3d)$, all $\mathbb Z/d$, and the positive even
groups through degree six vanish. There is no higher differential
in this range, by the base calculation. Universal coefficients,
starting in degree two, give
\[
 H_2=H_4=H_6=\mathbb Z/d,\qquad H_3=H_5=0.
\]
With coefficients $\mathbb Z/d$ the indicated transgressions vanish,
and $t,t^2,t^3$ survive; in particular the cube evaluates as a
generator of the dual of $H_6$. The same reasoning applies after
localization at every prime dividing $d$. At all other primes,
$K(\mathbb Z/d,2)$ is homologically trivial after localization,
as follows from its single nonzero homotopy group and the localized
Hurewicz theorem. This rules out additional two- or three-primary
classes omitted by passing to $R$.

Finally we use the product decomposition
\[
 K(A,2)=(\mathbb CP^\infty)^r\times K(\mathbb Z/d,2).
\]
The first factor has torsion-free homology concentrated in even
degrees, so the integral Kunneth formula has no Tor term here.
Its degree-six free part has rank $\binom{r+2}{3}$; tensoring its
degree-four, degree-two and degree-zero groups with the three
cyclic groups just computed gives respectively
$\binom{r+1}{2}$, $r$ and one cyclic summands. The stated products
evaluate on these distinct summands, and each is a perfect
detector in its coordinate. This proves the lemma.
\end{proof}
